\documentclass[11pt,a4paper]{article}

\usepackage[margin=1in]{geometry}
\usepackage{amsmath,amssymb}
\usepackage{booktabs}
\usepackage{array}
\usepackage{tabularx}
\usepackage{enumitem}
\usepackage{hyperref}
\usepackage[numbers,sort&compress]{natbib}

\hypersetup{
  colorlinks=true,
  linkcolor=blue,
  citecolor=blue,
  urlcolor=blue
}

\setlength{\parskip}{0.6em}
\setlength{\parindent}{0pt}
\renewcommand{\arraystretch}{1.15}
\urlstyle{same}

\title{Risk-Guided Hybrid Refinement cho tái tạo 3D từ ảnh UAV: nghiên cứu thực nghiệm trên Dronescapes và ODMData}
\author{Bản thảo paper thực nghiệm cho dự án GEO}
\date{Ngày 12 tháng 05 năm 2026}

\begin{document}
\maketitle

\begin{abstract}
Bài báo này trình bày một hướng triển khai thực dụng cho bài toán tái tạo hình học từ ảnh UAV, trong đó mục tiêu không phải là tối đa hóa chất lượng tuyệt đối bằng một pipeline đắt đỏ trên toàn bộ chuỗi ảnh, mà là \textbf{phân bổ chọn lọc ngân sách tinh chỉnh} vào những khung hình có rủi ro tái tạo cao. Phương pháp đề xuất, gọi là \textbf{risk-guided hybrid refinement}, sử dụng \textbf{DUSt3R} để suy ra hình học thô trên toàn bộ tập ảnh, chấm điểm rủi ro cho từng khung hình dựa trên thiếu hỗ trợ láng giềng, độ nghèo texture, thành phần depth prior và vị trí nhìn, sau đó chỉ chạy \textbf{MASt3R} trên tập con khung hình có rủi ro cao nhất. Phương pháp được đánh giá trên hai nguồn dữ liệu UAV thật: \textbf{Dronescapes} và \textbf{ODMData}. Kết quả cho thấy trên Dronescapes, phương pháp đề xuất đạt \textbf{RMSE 447.81} so với \textbf{447.94} của DUSt3R, đồng thời chỉ tinh chỉnh \textbf{25\%} số khung hình; trên ODMData ở chế độ pseudo-reference, phương pháp đề xuất cải thiện rõ rệt so với DUSt3R và tiến gần hơn tới lời giải tham chiếu. Tuy nhiên, kết quả cũng cho thấy phiên bản hiện tại \textbf{chưa vượt MASt3R full-pass về độ chính xác tuyệt đối} và \textbf{chưa thắng về thời gian chạy}, do vẫn chịu overhead của coarse pass, selective refine và khâu hợp nhất. Kết luận chính của bài báo là: selective refinement theo rủi ro là một hướng \textbf{có tín hiệu nghiên cứu tốt}, nhưng để trở thành một giải pháp vượt trội toàn diện thì còn cần tối ưu sâu hơn ở cấp độ mô hình và lịch lập lịch tính toán.
\end{abstract}

\section{Giới thiệu}

Tái tạo bản đồ 2D/2.5D/3D từ ảnh UAV là một bài toán trung tâm trong photogrammetry, bản đồ số và robot thị giác. Trong hơn một thập kỷ qua, cộng đồng đã phát triển rất mạnh các pipeline \textit{Structure-from-Motion} (SfM), \textit{Multi-View Stereo} (MVS) và các hệ visual mapping quy mô lớn cho ảnh UAV \citep{Colomina2014Unmanned,Jiang2017Efficient,Jiang2020Efficient,Liu2023Deep,Jiang2024Efficient}. Các phương pháp cổ điển như COLMAP cho phép phục hồi quỹ đạo camera và cấu trúc thưa khá ổn định \citep{Schonberger2016SfM}, trong khi các hệ MVS và neural reconstruction gần đây tiếp tục đẩy chất lượng hình học lên cao hơn nữa \citep{Liao2024High,Wu2025AerialBlocks}.

Tuy nhiên, bài toán thực tế lại không chỉ là ``dựng được hay không'', mà là \textbf{dựng với ngân sách tính toán nào}. Một pipeline đầy đủ kiểu SfM/MVS truyền thống có thể cho kết quả hình học tốt, nhưng chi phí chạy lớn và khó co giãn theo quy mô dữ liệu. Ở chiều ngược lại, các mô hình học sâu thế hệ mới như DUSt3R và MASt3R cho thấy tiềm năng rất cao cho tái tạo trực tiếp từ ảnh \citep{Wang2024DUSt3R,Leroy2024MASt3R,Wu2025AerialBlocks}, nhưng cũng đặt ra câu hỏi mới: có nhất thiết phải chạy mô hình nặng trên \textit{mọi} khung hình hay không?

Từ góc nhìn đó, bài báo này theo đuổi một luận điểm gọn nhưng thực dụng hơn: \textbf{không phải mọi khung hình đều đáng được tinh chỉnh như nhau}. Một số khung hình có thể đã đủ tốt sau bước tái tạo thô; một số khác lại mang rủi ro hình học cao do texture thấp, overlap yếu hoặc cấu hình góc nhìn bất lợi. Nếu có thể phát hiện đúng những vùng khó và chỉ tập trung ngân sách tính toán vào đó, pipeline có thể đạt điểm cân bằng tốt hơn giữa chất lượng và chi phí.

Trên cơ sở đó, bài báo đề xuất \textbf{risk-guided hybrid refinement}, một pipeline hai tầng:
\begin{enumerate}[leftmargin=1.5em]
  \item chạy DUSt3R trên toàn bộ tập ảnh để tạo ra một lớp hình học thô;
  \item chấm điểm rủi ro cho từng khung hình;
  \item chỉ refine các khung hình rủi ro cao bằng MASt3R;
  \item hợp nhất kết quả coarse và refine để thu được đầu ra cuối.
\end{enumerate}

Đóng góp chính của bài báo gồm bốn điểm:
\begin{enumerate}[leftmargin=1.5em]
  \item xây dựng một formulation \textbf{selective refinement theo rủi ro} cho tái tạo UAV;
  \item hiện thực hóa formulation này thành một pipeline benchmark khả dụng trong dự án GEO;
  \item đánh giá hệ thống trên \textbf{hai dataset thật} là Dronescapes và ODMData;
  \item rút ra một kết luận trung thực: phương pháp đề xuất \textbf{có ưu thế về trade-off chất lượng/chi phí}, nhưng \textbf{chưa phải phương án tốt nhất tuyệt đối}.
\end{enumerate}

\section{Công trình liên quan}

\subsection{Photogrammetry và MVS cho ảnh UAV}

Ảnh UAV đã trở thành nguồn dữ liệu quan trọng cho bản đồ 3D và khảo sát địa không gian \citep{Colomina2014Unmanned}. Các nghiên cứu về SfM hiệu quả cho ảnh UAV, đặc biệt ở ảnh oblique và quy mô lớn, đã cải thiện đáng kể lựa chọn cặp ảnh, hợp nhất graph và khả năng mở rộng \citep{Jiang2017Efficient,Jiang2020Efficient,Jiang2022Leveraging,Jiang2022Parallel,Jiang2024Efficient}. Song song đó, MVS dựa trên học sâu và các chiến lược nâng completeness cho urban scenes tiếp tục nâng chất lượng hình học, đặc biệt ở các khu vực phức tạp \citep{Liu2023Deep,Liao2024High}.

Tuy nhiên, các pipeline này nhìn chung vẫn có một đặc điểm chung: \textbf{chi ngân sách tái tạo tương đối đồng đều trên toàn bộ tập ảnh}. Đây là một giả định hợp lý ở pipeline offline chuẩn, nhưng lại kém linh hoạt nếu mục tiêu là giảm chi phí mà vẫn duy trì chất lượng đủ cao.

\subsection{Neural reconstruction và các mô hình pointmap}

Các mô hình pointmap gần đây đã thay đổi đáng kể bức tranh tái tạo hình học. DUSt3R đề xuất dự báo trực tiếp pointmaps và nới lỏng các ràng buộc projective truyền thống, từ đó đơn giản hóa đáng kể bài toán dựng hình học từ một hoặc nhiều ảnh \citep{Wang2024DUSt3R}. Trên nền đó, MASt3R mở rộng sang bài toán matching trong không gian 3D và cải thiện mạnh về độ mạnh matching cũng như downstream geometry \citep{Leroy2024MASt3R}.

Khi áp dụng lên aerial blocks, các mô hình này cho thấy tiềm năng lớn, đặc biệt trong các thiết lập góc nhìn khó hoặc sparse reconstruction \citep{Wu2025AerialBlocks}. Dù vậy, việc chạy full-pass các mô hình này trên mọi khung hình vẫn tốn kém, và câu hỏi về \textbf{phân bổ chọn lọc tính toán} vẫn còn khá mở.

\subsection{Uncertainty, confidence và selective processing}

Trong dense reconstruction, confidence và uncertainty từ lâu đã được xem là chìa khóa để kiểm soát lỗi \citep{Li2020Confidence}. Ở UAV mapping và robot mapping, nhiều hệ thống cũng đã tận dụng các chỉ báo độ tin cậy để quyết định nơi cần đo dày hơn hoặc nơi nên cập nhật world model mạnh hơn \citep{Campos2021ORB,Qin2018VINS,Wang2024Benchmarking}. Tuy nhiên, trên bài toán tái tạo UAV từ ảnh trong thiết lập hiện tại, vẫn còn thiếu một pipeline rõ ràng kết nối:
\begin{enumerate}[leftmargin=1.5em]
  \item một coarse pass phủ toàn cục,
  \item một risk model ở mức frame,
  \item và một refinement pass có ngân sách hạn chế.
\end{enumerate}

\subsection{Dataset và bối cảnh benchmark}

Để đánh giá ý tưởng theo hướng thực dụng, bài báo này dùng hai dataset thật có vai trò bổ sung nhau:
\begin{enumerate}[leftmargin=1.5em]
  \item \textbf{Dronescapes}, một benchmark UAV hướng semantic/depth/normals, hữu ích cho đo định lượng depth với ground-truth hoặc depth chuẩn hóa theo benchmark \citep{DronescapesDataset};
  \item \textbf{ODMData}, tập hợp các block ảnh photogrammetry thật do OpenDroneMap công bố để benchmark pipeline dựng bản đồ \citep{ODMDataBenchmarks}.
\end{enumerate}

Ngoài ra, bài báo cũng ý thức rằng các nhánh như 3D Gaussian Splatting đang rất mạnh ở novel-view synthesis và scene representation \citep{Yao2026ARSGaussian,Gao2025Enhanced}. Tuy nhiên, vì mục tiêu của nghiên cứu này là \textbf{depth-centric reconstruction under computational budget}, các phương pháp đó không được đưa vào benchmark chính của bài báo.

\section{Phương pháp đề xuất}

\subsection{Phát biểu bài toán}

Cho một tập ảnh UAV $\mathcal{I}=\{I_i\}_{i=1}^{N}$, mục tiêu là sinh ra một biểu diễn hình học đủ dùng cho mapping với chi phí nhỏ hơn so với việc refine toàn bộ tập ảnh bằng mô hình mạnh nhất. Ta muốn thiết kế một hàm chọn lọc sao cho chỉ một tập con $\mathcal{S}\subseteq \mathcal{I}$ được đưa qua bước refinement nặng, trong khi phần còn lại chỉ dùng kết quả coarse.

Ở đây, bài toán được viết lại thành:
\[
\min_{\mathcal{S}} \; \mathcal{L}_{geo}(\hat{\mathcal{D}}(\mathcal{I}, \mathcal{S})) + \lambda \, \mathcal{C}(\mathcal{S}),
\]
trong đó $\hat{\mathcal{D}}$ là depth/reconstruction cuối cùng sau khi hợp nhất coarse và refine, còn $\mathcal{C}(\mathcal{S})$ là chi phí tính toán tỷ lệ với số frame được refine.

\subsection{Kiến trúc pipeline}

Pipeline đề xuất gồm bốn bước:

\textbf{Bước 1: coarse reconstruction bằng DUSt3R.}
Toàn bộ tập ảnh được đưa qua DUSt3R để sinh ra depth hoặc pointmap thô. Bước này cung cấp lớp hình học nền trên toàn bộ scene.

\textbf{Bước 2: ước lượng rủi ro theo frame.}
Với mỗi frame, tính bốn thành phần:
\begin{enumerate}[leftmargin=1.5em]
  \item \textbf{overlap term} phản ánh mức thiếu hỗ trợ từ các ảnh láng giềng;
  \item \textbf{texture term} phản ánh độ nghèo texture cục bộ;
  \item \textbf{depth term} phản ánh độ bất định của depth prior nếu có;
  \item \textbf{view term} phản ánh mức bất lợi do vị trí nhìn trong chuỗi.
\end{enumerate}

Trong implementation hiện tại, điểm rủi ro được tính bởi:
\[
R_i = 0.35\,O_i + 0.25\,T_i + 0.20\,D_i + 0.20\,V_i,
\]
với $O_i, T_i, D_i, V_i \in [0,1]$ sau chuẩn hóa min-max. Cụ thể:
\begin{itemize}[leftmargin=1.5em]
  \item $O_i$ tăng khi cosine similarity với các frame láng giềng giảm;
  \item $T_i$ tăng khi độ lệch chuẩn mức xám giảm, tức texture nghèo hơn;
  \item $D_i$ tăng khi depth prior có độ phân tán lớn; trong benchmark hiện tại, thành phần này bị vô hiệu vì không bật prior depth ngoại sinh;
  \item $V_i$ tăng khi frame ở xa vùng trung tâm của chuỗi.
\end{itemize}

\textbf{Bước 3: selective refinement bằng MASt3R.}
Chỉ các frame có điểm $R_i$ cao nhất mới được refine bằng MASt3R. Trong benchmark hiện tại, ngân sách refine được đặt là:
\[
K = 8
\]
frame trên mỗi dataset.

\textbf{Bước 4: merge coarse và refine.}
Depth đã refine sẽ ghi đè hoặc hợp nhất với coarse depth trên tập frame được chọn, tạo ra đầu ra cuối cho benchmark.

\subsection{Bản chất của novelty}

Điểm mới của phương pháp không nằm ở việc phát minh một backbone 3D mới, mà ở cách \textbf{biến quyết định ``refine ở đâu'' thành đối tượng tối ưu chính}. Nói cách khác, bài báo đặt trọng tâm vào \textbf{tính kinh tế của refinement}, chứ không chỉ vào chất lượng cuối.

\section{Thiết lập thực nghiệm}

\subsection{Dataset}

\textbf{ODMData.} Bài báo dùng sample \texttt{mygla} từ bộ benchmark chính thức của OpenDroneMap \citep{ODMDataBenchmarks}. Sau bước chuẩn bị, dataset này có \textbf{29 frame} trong cấu hình benchmark hiện tại.

\textbf{Dronescapes.} Bài báo dùng split \texttt{test\_set\_annotated\_only} của Dronescapes \citep{DronescapesDataset}, và lấy \textbf{32 frame} cho bài toán POC. Một chi tiết quan trọng về implementation là phần \texttt{rgb} trên nguồn dữ liệu Hugging Face của Dronescapes được lưu dưới dạng \texttt{.npz}; pipeline của dự án đã chuyển đổi chúng về \texttt{.png} ở bước export để phục vụ benchmark thống nhất.

\subsection{Phương pháp so sánh}

Ba phương pháp được đưa vào thực nghiệm hiện tại là:
\begin{enumerate}[leftmargin=1.5em]
  \item \textbf{DUSt3R full-pass} (\texttt{dust3r\_real}): chạy DUSt3R trên toàn bộ tập ảnh;
  \item \textbf{MASt3R full-pass} (\texttt{mast3r\_real}): chạy MASt3R trên toàn bộ tập ảnh;
  \item \textbf{risk-guided hybrid refinement} (\texttt{risk\_hybrid\_real}): coarse DUSt3R + chọn top-$K$ frame theo risk + refine bằng MASt3R.
\end{enumerate}

Framework GEO cũng hỗ trợ nhánh \texttt{COLMAP + OpenMVS}, nhưng run thực nghiệm được báo cáo trong bài này là một run POC neural-centric nên không bao gồm baseline photogrammetry cổ điển.

\subsection{Thiết lập tham số}

Run thực nghiệm trên máy \texttt{avis@100.96.109.77} sử dụng:
\begin{enumerate}[leftmargin=1.5em]
  \item \textbf{coarse image size}: 224;
  \item \textbf{refine image size}: 384;
  \item \textbf{batch size}: 1;
  \item \textbf{window size}: 2;
  \item \textbf{top-$K$ refine frames}: 8;
  \item \textbf{risk neighbors}: 4;
  \item \textbf{coarse/refine device}: CUDA.
\end{enumerate}

Phần cứng chạy benchmark gồm:
\begin{enumerate}[leftmargin=1.5em]
  \item CPU Intel Core Ultra 9 285K;
  \item RAM 60 GiB;
  \item GPU NVIDIA RTX 5000 Ada 32 GiB;
  \item Ubuntu 24.04.4.
\end{enumerate}

\subsection{Metric đánh giá}

Trên Dronescapes, bài báo dùng:
\begin{enumerate}[leftmargin=1.5em]
  \item \textbf{MAE},
  \item \textbf{RMSE},
  \item \textbf{AbsRel},
  \item \textbf{valid ratio},
  \item \textbf{runtime}.
\end{enumerate}

Trên ODMData, vì không có ground-truth depth trực tiếp cho thiết lập benchmark hiện tại, hệ thống dùng \textbf{pseudo-reference evaluation}. Ở run này, \texttt{MASt3R} được chọn làm \textbf{reference method}, nên sai số của chính \texttt{mast3r\_real} trên ODMData bằng 0 theo định nghĩa và chỉ nên được xem là mốc neo để so sánh khoảng cách của các phương pháp khác.

\section{Kết quả}

\subsection{Kết quả trên Dronescapes}

\begin{table}[t]
\centering
\small
\begin{tabularx}{\textwidth}{>{\raggedright\arraybackslash}p{3.1cm} >{\centering\arraybackslash}p{1.8cm} >{\centering\arraybackslash}p{1.8cm} >{\centering\arraybackslash}p{1.6cm} >{\centering\arraybackslash}p{1.8cm} >{\centering\arraybackslash}p{1.7cm}}
\toprule
\textbf{Phương pháp} & \textbf{MAE} & \textbf{RMSE} & \textbf{AbsRel} & \textbf{Runtime (s)} & \textbf{Selected ratio} \\
\midrule
DUSt3R & 345.054 & 447.936 & 0.9905 & 182.67 & 1.00 \\
MASt3R & \textbf{343.628} & \textbf{446.506} & \textbf{0.9858} & 189.98 & 1.00 \\
Risk-hybrid & 344.913 & 447.812 & 0.9900 & 363.64 & 0.25 \\
\bottomrule
\end{tabularx}
\caption{Kết quả trên \texttt{dronescapes\_poc\_real}.}
\label{tab:dronescapes-results}
\end{table}

Bảng \ref{tab:dronescapes-results} cho thấy:
\begin{enumerate}[leftmargin=1.5em]
  \item MASt3R full-pass đạt chất lượng tốt nhất về mọi metric chính;
  \item risk-hybrid cải thiện nhẹ so với DUSt3R ở cả MAE, RMSE và AbsRel;
  \item risk-hybrid chỉ refine \textbf{25\%} số frame, nhưng vẫn tiến khá sát MASt3R.
\end{enumerate}

Đây là tín hiệu tốt cho luận điểm selective refinement. Tuy nhiên, runtime của risk-hybrid lại cao hơn đáng kể, do pipeline hiện tại vẫn phải trả giá cho cả:
\begin{enumerate}[leftmargin=1.5em]
  \item coarse pass toàn cục bằng DUSt3R,
  \item refine pass cục bộ bằng MASt3R,
  \item và bước merge hậu xử lý.
\end{enumerate}

\subsection{Kết quả trên ODMData}

\begin{table}[t]
\centering
\small
\begin{tabularx}{\textwidth}{>{\raggedright\arraybackslash}p{3.1cm} >{\centering\arraybackslash}p{1.8cm} >{\centering\arraybackslash}p{1.8cm} >{\centering\arraybackslash}p{1.6cm} >{\centering\arraybackslash}p{1.8cm} >{\centering\arraybackslash}p{1.7cm}}
\toprule
\textbf{Phương pháp} & \textbf{MAE} & \textbf{RMSE} & \textbf{AbsRel} & \textbf{Runtime (s)} & \textbf{Selected ratio} \\
\midrule
DUSt3R & 0.1601 & 0.2492 & 0.0991 & 196.92 & 1.00 \\
MASt3R (reference) & 0.0000 & 0.0000 & 0.0000 & 207.64 & 1.00 \\
Risk-hybrid & \textbf{0.1202} & \textbf{0.2143} & \textbf{0.0772} & 456.44 & 0.276 \\
\bottomrule
\end{tabularx}
\caption{Kết quả trên \texttt{odmdata\_odm\_poc\_mygla\_real}. Sai số của MASt3R bằng 0 vì đây là pseudo-reference cho run này.}
\label{tab:odm-results}
\end{table}

Trên ODMData, cần đọc Bảng \ref{tab:odm-results} cẩn thận. Vì MASt3R là pseudo-reference của run, hàng \texttt{mast3r\_real} không phải là một chiến thắng tuyệt đối mà chỉ là mốc neo. Điều đáng chú ý là:
\begin{enumerate}[leftmargin=1.5em]
  \item risk-hybrid tốt hơn rõ rệt so với DUSt3R trên mọi metric;
  \item risk-hybrid chỉ refine khoảng \textbf{27.6\%} số frame;
  \item khoảng cách của risk-hybrid tới reference nhỏ hơn đáng kể so với DUSt3R.
\end{enumerate}

Điều này cho thấy, trong không gian không có ground truth tuyệt đối nhưng có một lời giải tham chiếu hợp lý, selective refinement thực sự giúp kéo coarse reconstruction lại gần lời giải mạnh hơn.

\section{Thảo luận}

\subsection{Risk-hybrid hiện tốt ở đâu}

Kết quả hiện tại cho thấy \texttt{risk\_hybrid\_real} có ba ưu điểm rõ:
\begin{enumerate}[leftmargin=1.5em]
  \item \textbf{ổn định hơn DUSt3R} về chất lượng hình học trên cả hai dataset;
  \item \textbf{không cần refine toàn bộ tập ảnh};
  \item \textbf{tiến gần MASt3R} trong khi chỉ dùng một phần ngân sách refine.
\end{enumerate}

Nói cách khác, điểm mạnh thực sự của phương pháp nằm ở \textbf{hiệu quả sử dụng refinement budget}, không nằm ở việc giành best absolute accuracy.

\subsection{Risk-hybrid hiện chưa tốt ở đâu}

Điểm yếu lớn nhất của pipeline hiện tại là:
\begin{enumerate}[leftmargin=1.5em]
  \item \textbf{runtime chưa tốt}. Trên cả Dronescapes và ODMData, risk-hybrid đều chậm hơn MASt3R full-pass trong run hiện tại;
  \item \textbf{risk score còn heuristic}. Công thức hiện tại dựa trên overlap, texture, view và depth prior đơn giản, chưa học được uncertainty thực sự từ backbone;
  \item \textbf{refinement ở mức frame còn thô}. Có thể nhiều scene chỉ cần refine vài tile hoặc patch nhỏ thay vì cả frame;
  \item \textbf{đánh giá ODMData còn là pseudo-reference}. Điều này hữu ích cho development, nhưng không đủ để đưa ra kết luận tuyệt đối về hình học.
\end{enumerate}

\subsection{Ý nghĩa nghiên cứu}

Từ góc nhìn học thuật, bài báo này chưa đủ để kết luận rằng selective refinement là giải pháp mạnh nhất. Nhưng nó đưa ra một thông điệp quan trọng:
\begin{quote}
\textit{Selective refinement theo rủi ro có thể tạo ra chất lượng gần với full refinement, ngay cả khi chỉ tinh chỉnh một phần nhỏ số frame.}
\end{quote}

Nếu được tối ưu tốt hơn ở cấp độ hệ thống, đây là một hướng rất có triển vọng cho UAV reconstruction under budget.

\subsection{Hướng phát triển tiếp theo}

Bốn hướng tiếp theo hợp lý nhất là:
\begin{enumerate}[leftmargin=1.5em]
  \item học một \textbf{risk predictor} tốt hơn thay cho heuristic hiện tại;
  \item chuyển từ \textbf{frame-level refinement} sang \textbf{tile-level refinement};
  \item tái sử dụng feature hoặc latent của coarse pass để giảm overhead của refine pass;
  \item đưa \texttt{COLMAP + OpenMVS} trở lại benchmark khi hạ tầng phần cứng ổn định hơn, nhằm hoàn thiện so sánh với nhánh photogrammetry cổ điển.
\end{enumerate}

\section{Kết luận}

Bài báo này trình bày một phiên bản thực dụng của bài toán tái tạo 3D từ ảnh UAV, trong đó trọng tâm không phải là dựng hình học mạnh nhất có thể bằng full-pass trên mọi khung hình, mà là \textbf{chọn đúng nơi để chi refinement budget}. Trên hai dataset thật là Dronescapes và ODMData, phương pháp \texttt{risk\_hybrid\_real} cho thấy khả năng cải thiện nhất quán so với DUSt3R coarse-only và tiến gần kết quả của MASt3R, trong khi chỉ refine một phần nhỏ số frame.

Tuy vậy, kết quả hiện tại cũng chỉ ra giới hạn quan trọng: phiên bản implementation hiện tại \textbf{chưa vượt MASt3R về chất lượng tuyệt đối} và \textbf{chưa đạt lợi thế runtime}. Vì vậy, đóng góp của bài báo nên được hiểu đúng ở mức: \textbf{một bằng chứng thực nghiệm rằng selective refinement theo rủi ro là hướng nghiên cứu hợp lý}, chứ chưa phải một lời tuyên bố về một phương pháp tốt nhất tuyệt đối.

Nhìn từ góc độ dự án GEO, kết quả này là đủ để khẳng định rằng hướng nghiên cứu hiện tại đáng để tiếp tục đầu tư. Giá trị chính của hệ thống không còn nằm ở một benchmark đơn thuần, mà ở việc biến bài toán ``reconstruct everything'' thành ``reconstruct selectively and intelligently''.

\nocite{Zhao2014Direct,Zheng2018multi,Wu2018Integration,Zhang2020UAV,Yu2021Automatic,Xu2021Robust,Chen2024End,Tong2024Large,Wu2024evaluation,Shi2025Accurate}

\bibliographystyle{plainnat}
\bibliography{references}

\end{document}
